\documentclass[a4paper]{article}

\usepackage[UKenglish]{babel}
\usepackage{graphicx}
\usepackage{amsmath}
\usepackage{amssymb}
\usepackage[T1]{fontenc}
\usepackage{listings}

\AtBeginDocument{\renewcommand{\abstractname}{Abstract}}

\begin{document}
\begin{titlepage}
	\begin{center}
	\textsc{\LARGE Machine Learning in Practice\\}
	\textsc{\Large Report}\\[1.5cm]
	\includegraphics[height=100pt]{logo}
   
	\vspace{0.4cm}
	\textsc{\Large Radboud University Nijmegen}\\[.5cm]
	\hrule
	\vspace{0.4cm}
	\textbf{\huge Acquire Valued Shoppers Challenge}\\[0.4cm]
	\hrule
	\vspace{2cm}
	\begin{minipage}[t]{0.45\textwidth}
	\begin{flushleft} \large
	Wouter Geraedts\\
	s0814857\\[0.7cm]
	Matthijs Hendriks\\
	s4068459\\[0.7cm]
	\end{flushleft}
	\end{minipage}
	\begin{minipage}[t]{0.45\textwidth}
	\begin{flushright} \large
	Thomas N\"agele\\
	sXXXXXX\\[0.7cm]
	Mathijs Vos\\
	s4024443\\[0.7cm]
	\end{flushright}
	\end{minipage}
	\vspace{.7cm}
	
	\begin{abstract}
		ABSTRACT HIERZO
	\end{abstract}
	\vspace{.7cm}

	{\large 3 July 2014}
	\vfill
	\end{center}

\end{titlepage}

\newpage

\section{Introduction}


\section{Approach}

The Acquire Valued Shoppers Challenge requires a very specific kind of machine
learning due to the large dataset. Loading the entire set in memory requires
specialized equipment, which we had access to. This approach makes analyzing 
this dataset not particulary convenient.

This challenge was approached in a typical manner: first a basic pipeline was
implemented in order to enable rapid incremental improvements. The trainingset
was split in three subsets (train, test and tuning) in order to measure the
classification performance. Then we could add, analyze and prune new features.
Also a few classification algorithms were tried and their parameters tuned.

The \texttt{shopper-hivemind} application features a streaming pipeline for
converting the raw CSV dataset into relational objects. These objects are
subsequently fed to the feature extractor. A feature matrix is generated,
normalized and then fitted by a classifier. This classifier can then immediately
be used to predict test data.

\subsection{Dataset loading}

The original dataset is 30GiB of raw CSV. Because loading this dataset took a
quite some time when implemented in Python, we decided to use C++ instead.
Loading this data is implemented as a streaming set of readers. This means that
during this phase only a single data row is in memory.

\begin{itemize}
    \item The Gzipped CSV file is read to a simple datarecord.
    \item These datarecords are cached in an intermediate bytepacked file, in the
        \emph{Msgpack}\footnote{\url{http://msgpack.org/}} format (also Gzipped).
    \item Subsequently these datarecords are linked together to form complex
        relational dataobjects, and are again cached in a Gzipped Msgpack file.
        These files use roughly 2.4GiB of diskspace.
        % Graph of datastructures (?)
    \item Finally feature extraction is performed on each object.
\end{itemize}

These cache files are stored in and read directly from a Gzipped state. Without
compression the disk throughput bottlenecked the loading of the dataset. The
required bandwidth even outclassed the performance of some consumergrade SSD's.
This approach solves that problem, at the cost of some CPU time.

Most classification algorithms require the features to be presented as an
in-memory feature matrix. Luckily our feature matrix is only 1.4GiB, which fits
in the memory of most modern machines.

Finally a method for creating uniformly distributed subsets of the dataset was
implemented. We used a single split for training, testing and tuning. Also a
1000 row subset was created for debugging and rapid prototyping.

In practice we did not use cross validation actively, nor did we measure our
performance using this tuning and testset. Because fitting the models took quite
some time, we only measured our performance by predicting the final test set and
submitting the entry to Kaggle.

\section{Features}


\section{Classification}


\section{Results}


\section{Conclusions}


\end{document}
